
\medskip
\phantomsection\label{lem:sga3-VIB-11.12}%
\noindent\textbf{Lemma 11.12.---}
\emph{Let \(k\) be a field and let \(G\) be an affine
\(k\)-group. Put \(A=\mathcal A(G)\). Given \(x\in A\),
there is a finitely generated \(k\)-subalgebra \(B\) of
\(A\) such that}
\[
  x\in B,\qquad
  \Delta(B)\subset B\otimes_kB,\qquad
  u(B)\subset B,
\]
\emph{where \(u\) denotes the involution of \(A\)
corresponding to the inversion morphism of \(G\).}%
\footnote{\label{note:sga3-VIB-125}Editor's note: on the
one hand, this is generalized in \hyperref[sec:sga3-VIB-13]{Section 13} to the case
where \(G\) is affine and flat over a regular base of
dimension at most \(2\). On the other hand, the original
has been expanded and corrected in what follows.}

We may suppose that \(x\neq0\). Then
\(\Delta(x)\neq0\), since
\[
  (\varepsilon\otimes\mathrm{id})\Delta(x)=x,
\]
where \(\varepsilon\) denotes the augmentation
\textup{(counit)} of \(A\). Write
\phantomsection\label{form:sga3-VIB-11.12-eq1}%
\[
  \tag{1}
  \Delta(x)=\sum_{j=1}^{n}e_j\otimes a_j,
  \qquad n\ \text{minimal}.
\]
The \(e_j\) \textup{(respectively the \(a_j\))} are then
linearly independent. Complete
\((e_1,\ldots,e_n)\) to a basis
\((e_i)_{i\in I}\) of \(A\), and, for
\[
  j\in J=\{1,\ldots,n\},
\]
write
\[
  \Delta(e_j)=\sum_{i\in I}e_i\otimes b_{ij}.
\]
Applying \(\Delta\otimes\mathrm{id}\) and
\(\mathrm{id}\otimes\Delta\) to
\hyperref[form:sga3-VIB-11.12-eq1]{\textup{(1)}}, the
axiom
\hyperref[ax:sga3-VIB-11.8.0-CO1]{\textup{(CO 1)}} of
\hyperref[def:sga3-VIB-11.8.0]{11.8.0}
\textup{(see also \textup{(HA 1)} in I, \SGARefLink{sga3:i:subsection:4-2}{4.2})} gives
\[
  \sum_{j\in J}e_j\otimes\Delta(a_j)
  =
  \sum_{\ell\in J}\Delta(e_\ell)\otimes a_\ell
  =
  \sum_{i\in I}e_i\otimes
    \left(\sum_{\ell\in J}b_{i\ell}\otimes a_\ell\right).
\]
Since the \(e_i\) are linearly independent, it follows that
\phantomsection\label{form:sga3-VIB-11.12-eq2}%
\[
  \tag{2}
  \Delta(a_j)=\sum_{\ell\in J}b_{j\ell}\otimes a_\ell
  \qquad (j\in J).
\]

Let \(B\) be the finitely generated \(k\)-subalgebra of
\(A\) generated by the \(b_{ij}\) and the \(u(b_{ij})\),
for \(i,j\in J\). It is clear that \(u(B)=B\).

Applying \(\Delta\otimes\mathrm{id}\) and
\(\mathrm{id}\otimes\Delta\) to
\hyperref[form:sga3-VIB-11.12-eq2]{\textup{(2)}}, we obtain
once again from \textup{(CO 1)} the equalities
\[
  \sum_{\ell\in J}\Delta(b_{j\ell})\otimes a_\ell
  =
  \sum_{i\in J}b_{ji}\otimes\Delta(a_i)
  =
  \sum_{i,\ell\in J}b_{ji}\otimes b_{i\ell}\otimes a_\ell.
\]
Since the \(a_\ell\) are linearly independent, this yields
\phantomsection\label{form:sga3-VIB-11.12-eq3}%
\[
  \tag{3}
  \Delta(b_{j\ell})
    =\sum_{i\in J}b_{ji}\otimes b_{i\ell}
  \qquad (j,\ell\in J).
\]
Moreover,
\[
  \Delta\circ u=(u\otimes u)\circ v\circ\Delta,
  \qquad v(a\otimes b)=b\otimes a,
\]
and hence
\phantomsection\label{form:sga3-VIB-11.12-eq4}%
\[
  \tag{4}
  \Delta\bigl(u(b_{j\ell})\bigr)
    =\sum_{i\in J}u(b_{i\ell})\otimes u(b_{ji})
  \qquad (j,\ell\in J).
\]
Since \(\Delta\) is an algebra homomorphism,
\hyperref[form:sga3-VIB-11.12-eq3]{\textup{(3)}} and
\hyperref[form:sga3-VIB-11.12-eq4]{\textup{(4)}} imply
\[
  \Delta(B)\subset B\otimes_kB.
\]
Finally, axiom
\hyperref[ax:sga3-VIB-11.8.0-CO2]{\textup{(CO 2)}} of
\hyperref[def:sga3-VIB-11.8.0]{11.8.0}
\textup{(see also \textup{(HA 2)} in I, \SGARefLink{sga3:i:subsection:4-2}{4.2})}, applied to
\hyperref[form:sga3-VIB-11.12-eq2]{\textup{(2)}} and
\hyperref[form:sga3-VIB-11.12-eq1]{\textup{(1)}}, gives%
\begingroup
\renewcommand{\thefootnote}{\fnsymbol{footnote}}%
\footnote[1]{\label{note:sga3-VIB-11.12-counit-source-reading}
Translator's note: the French re-edition prints
\(a_j=\sum_{i\in I}\varepsilon(a_j)b_{ij}\).
Applying \(\mathrm{id}\otimes\varepsilon\) to
\hyperref[form:sga3-VIB-11.12-eq2]{\textup{(2)}} gives the
first identity displayed here. Reinhold's English translation repeats the printed formula.}
\endgroup
\phantomsection\label{form:sga3-VIB-11.12-counit-identities}%
\[
  a_j=\sum_{i\in J}b_{ji}\varepsilon(a_i),
  \qquad
  x=\sum_{j\in J}\varepsilon(e_j)a_j.
\]
Thus \(x\in B\), proving the lemma.

\medskip
\phantomsection\label{prop:sga3-VIB-11.13}%
\noindent\textbf{Proposition 11.13.---}
\emph{Let \(k\) be a field and let \(G\) be an affine
\(k\)-group with algebra \(A\). Then \(G\) is the
projective limit of an increasing filtered system of
affine \(k\)-groups of finite type whose transition
morphisms are faithfully flat.}

If \(B\) and \(B'\) are two finitely generated
subalgebras of \(A\) stable under \(\Delta\) and \(u\), the
subalgebra generated by \(B\) and \(B'\) has the same
properties. Hence, by
\hyperref[lem:sga3-VIB-11.12]{Lemma 11.12}, \(A\) is the
direct limit of an increasing filtered family
\((B_i)_{i\in I}\) of finitely generated subalgebras stable
under \(\Delta\) and \(u\). Each \(B_i\), equipped with the
restriction of \(u\) and with the morphism
\[
  B_i\longrightarrow B_i\otimes_kB_i
\]
induced by \(\Delta\), is a Hopf algebra. Thus, by I, \SGARefLink{sga3:i:subsection:4-2}{4.2},
it is the algebra of an affine \(k\)-group \(G_i\) of
finite type over \(k\). Finally, since
\[
  A=\varinjlim_i B_i,
\]
we have
\[
  G=\varprojlim_i G_i
\]
\textup{(cf.\ EGA IV\(_3\), 8.2.3)}. The transition
morphisms are faithfully flat by the following lemma.

\medskip
\phantomsection\label{lem:sga3-VIB-11.14}%
\noindent\textbf{Lemma 11.14.---}
\emph{Let \(k\) be a field, let
\(u:G\to H\) be a morphism between affine \(k\)-groups of
finite type, and let}
\[
  u^\natural:B\longrightarrow A
\]
\emph{be the corresponding morphism of \(k\)-algebras.}%
\emph{Then \(u\) is faithfully flat if and only if
\(u^\natural\) is injective.}

The condition is plainly necessary
\textup{(cf.\ EGA 0\(_{\mathrm I}\), 6.6.1)}. Conversely,
suppose that \(u^\natural\) is injective, and put
\[
  N=\operatorname{Ker}(u).
\]
By
\hyperref[thm:sga3-VIA-3.3.2]{VI\(_A\), 3.3.2} and
\hyperref[prop:sga3-VIA-5.4.1]{VI\(_A\), 5.4.1},
\(G/N\) is a \(k\)-group of finite type and \(u\) factors
as
\[
  G\xrightarrow{\ p\ }G/N\xrightarrow{\ v\ }H,
\]
where \(p\) is faithfully flat and \(v\) is a closed
immersion. Since \(H\) is affine, \(G/N\) is affine, and
\[
  v^\natural:B\longrightarrow\mathcal O(G/N)
\]
is surjective \textup{(cf.\ EGA I, 4.2.3)}. But
\[
  u^\natural=p^\natural\circ v^\natural
\]
is injective, so \(v^\natural\) is also injective. Hence
\(v^\natural\), and therefore \(v\), is an isomorphism.
Since \(p\) is faithfully flat, so is \(u\).

\medskip
\phantomsection\label{def:sga3-VIB-11.15}%
\noindent\textbf{Definition 11.15.---}%
\footnote{\label{note:sga3-VIB-127}Editor's note: we have
added the hypothesis that \(G\) be quasi-compact and have
spelled out the equivalence of conditions
\textup{(i)} and \textup{\SGARefLink{concl:sga3-VIB-11.15-i}{(ii)}}.}
\emph{Let \(k\) be a field, let \(G\) be a quasi-compact
\(k\)-group, and let \(V\) be a \(k\)-vector space endowed
with a \(k\)-linear action of \(G\), hence, by
\hyperref[prop:sga3-VIB-11.6.1]{11.6.1}, with an
\(\mathcal A(G)\)-comodule structure}
\[
  \delta:V\longrightarrow V\otimes\mathcal A(G).
\]
\emph{Let \(v\in V\) be nonzero. The following conditions
are equivalent:}
\begin{enumerate}
\item[\textup{(i)}]\phantomsection\label{concl:sga3-VIB-11.15-i}
  \emph{There is a \(\lambda\in\mathcal A(G)\),
  necessarily unique, such that}
  \[
    \delta(v)=v\otimes\lambda.
  \]

\item[\textup{(ii)}]\phantomsection\label{concl:sga3-VIB-11.15-ii}
  \emph{For every \(k\)-algebra \(R\) and every
  \(g\in G(R)\), one has \(g\cdot v\in Rv\); that is,
  there is a necessarily unique \(f\in R\) such that,
  in \(V\otimes R\),}
  \[
    g\cdot v=v\otimes f.
  \]
\end{enumerate}

It is clear that
\hyperref[concl:sga3-VIB-11.15-i]{\textup{(i)}}
implies
\hyperref[concl:sga3-VIB-11.15-ii]{\textup{(ii)}}.
Conversely, if \textup{\SGARefLink{concl:sga3-VIB-11.15-i}{(ii)}} holds, apply it with
\[
  R=\mathcal A(G),
  \qquad
  g=\operatorname{id}_{\mathcal A(G)}.
\]
There is then a unique \(\lambda\in\mathcal A(G)\) such
that \(\delta(v)=v\otimes\lambda\).

When \(v\) satisfies these conditions, we call \(v\) a
\emph{semi-invariant vector under \(G\)}, and
\(\lambda\) the \emph{weight} of \(v\). We also say that
\(v\) is a \emph{semi-invariant of weight \(\lambda\)}.

Let \(\Delta\) denote the comultiplication of
\(\mathcal A(G)\). The equality
\phantomsection\label{form:sga3-VIB-11.15-group-like}%
\[
  v\otimes\lambda\otimes\lambda
  =
  (\delta\otimes\mathrm{id})(\delta(v))
  =
  (\mathrm{id}\otimes\Delta)(\delta(v))
  =
  v\otimes\Delta(\lambda)
\]
implies
\[
  \Delta(\lambda)=\lambda\otimes\lambda.
\]
Consequently, \(\lambda\) defines a morphism of Hopf
algebras
\[
  \mathcal A(\mathbb G_{m,k})
    =k[T,T^{-1}]
    \longrightarrow\mathcal A(G),
  \qquad T\longmapsto\lambda,
\]
and hence a morphism of \(k\)-groups
\[
  \lambda:G\longrightarrow\mathbb G_{m,k}.
\]
Thus \(\lambda\) is a \emph{character} of \(G\), called the
\emph{character associated with the semi-invariant vector
\(v\)}.

\medskip
\phantomsection\label{lem:sga3-VIB-11.16.0}%
\noindent\textbf{Lemma 11.16.0.---}%
\footnote{\label{note:sga3-VIB-128}Editor's note: this
lemma, taken from \textup{[DG70]}, Section II.2, \SGARefLink{prop:sga3-VIB-3.5}{3.5}, has
been added.}
\emph{Let \(k\) be a field, let \(H\) be an affine
\(k\)-group, let \(V\) be an \(H\)-module of dimension
\(n\), and let \(U\) be a \(d\)-dimensional vector
subspace of \(V\). Consider the line}
\[
  D=\bigwedge\nolimits^dU
    \subset\bigwedge\nolimits^dV.
\]
\emph{Then \(U\) is stable under \(H\) if and only if
\(D\) is stable under \(H\).}

Necessity is clear. To prove sufficiency, we may assume
\(d<n\). Let \((e_1,\ldots,e_d)\) be a basis of \(U\), and
complete it to a basis \((e_1,\ldots,e_n)\) of \(V\).
For every \(k\)-algebra \(R\), the module
\[
  V_R=V\otimes R
\]
is free over \(R\), and
\phantomsection\label{form:sga3-VIB-11.16.0-UR}%
\[
  U_R
  =
  \{\,v\in V_R\mid
    v\wedge(e_1\wedge\cdots\wedge e_d)=0\,\}.
\]
Indeed, for \(i>d\), the elements
\[
  e_i\wedge e_1\wedge\cdots\wedge e_d
\]
are linearly independent in
\(\bigwedge_R^{d+1}V_R\). Since \(H(R)\) acts on
\(\bigwedge_R^\bullet V_R\) by
\[
  h(x_1\wedge\cdots\wedge x_s)
  =
  h(x_1)\wedge\cdots\wedge h(x_s),
\]
it follows that, if \(H(R)\) stabilizes
\[
  D_R=R(e_1\wedge\cdots\wedge e_d),
\]
then it also stabilizes \(U_R\).

\medskip
\phantomsection\label{thm:sga3-VIB-11.16}%
\noindent\textbf{Theorem 11.16 \textup{(Chevalley)}.---}
\emph{Let \(k\) be a field, let \(G\) be an affine
algebraic \(k\)-group, and let \(H\) be a closed group
subscheme of \(G\).}%
\footnote{\label{note:sga3-VIB-129}Editor's note: first,
recall that every group subscheme of \(G\) is closed
\textup{(\hyperref[cor:sga3-VIB-1.4.2]{1.4.2})}. Second,
we have stated the result in its usual form, namely that
``\(H\) is the stabilizer of a line in a representation of
\(G\),'' while retaining the original formulation in terms
of a sequence \(a_1,\ldots,a_n\) of semi-invariants in
\(\mathcal A(G)\).}
\emph{Then there are a finite-dimensional \(G\)-module
\(V\) and a line \(D\) in \(V\) such that}
\[
  H=\operatorname{Norm}_G(D);
\]
\emph{that is, for every \(k\)-algebra \(R\),}
\phantomsection\label{form:sga3-VIB-11.16-stabilizer}%
\[
  H(R)=\{\,g\in G(R)\mid g(D_R)=D_R\,\}.
\]

\emph{Equivalently, there are finitely many elements
\(a_1,\ldots,a_n\in\mathcal A(G)\), all semi-invariant of
the same weight \(\lambda\) for the ``right'' action of
\(H\); that is,}
\[
  a_i(gh)=\lambda(h)a_i(g)
\]
\emph{for every \(k\)-algebra \(R\),
\(g\in G(R)\), and \(h\in H(R)\), such that \(H\) is the
largest closed group subscheme of \(G\) under which the
\(a_i\) are semi-invariant.}

\emph{Proof.}
Let \(\Delta\) and \(\varepsilon\) denote, respectively,
the comultiplication and the augmentation of
\[
  A=\mathcal A(G).
\]
There is an ideal \(I\) of \(A\), contained in
\(\operatorname{Ker}\varepsilon\), such that
\[
  H=\operatorname{Spec}(A/I),
  \qquad
  \Delta(I)\subset I\otimes A+A\otimes I.
\]
Put \(B=A/I\), and let \(\pi:A\to B\) be the projection.
Consider the left action of \(H\) on \(A\) given by
\[
  (h\phi)(g)=\phi(gh).
\]
The corresponding \(B\)-comodule structure is
\phantomsection\label{form:sga3-VIB-11.16-H-coaction-A}%
\[
  \overline{\Delta}:\quad
  A\xrightarrow{\ \Delta\ }A\otimes A
    \xrightarrow{\ \mathrm{id}_A\otimes\pi\ }A\otimes B.
\]
Then \(I\) is an \(H\)-submodule of \(A\), since
\[
  \overline{\Delta}(I)\subset I\otimes B.
\]

The \(k\)-algebra \(A\) is finitely generated, and hence
Noetherian, so \(I\) has a finite system of generators
\((x_1,\ldots,x_r)\). By
\hyperref[lem:sga3-VIB-11.8]{11.8}, the \(x_i\) are
contained in a finite-dimensional \(G\)-submodule \(V\) of
\(A\). Then
\[
  W=V\cap I
\]
is a finite-dimensional \(H\)-module; write
\[
  d=\dim_kW.
\]
Since \(V\) contains every \(x_i\), \(W\) generates the
ideal \(I\).

Put
\[
  E=\bigwedge\nolimits^dV,
\]
choose a basis \((w_1,\ldots,w_d)\) of \(W\), and choose a
basis \((e_0,\ldots,e_n)\) of \(E\) containing
\[
  e_0=w_1\wedge\cdots\wedge w_d.
\]
The action of \(G\) on \(V\) canonically induces an action
of \(G\) on \(E\), hence an \(A\)-comodule structure
\[
  \rho:E\longrightarrow E\otimes A.
\]
For \(j=0,\ldots,n\), write%
\begingroup
\renewcommand{\thefootnote}{\fnsymbol{footnote}}%
\footnote[2]{\label{note:sga3-VIB-11.16-rho-index-source-reading}
Translator's note: the French re-edition and Reinhold's English translation print the lower summation bound as
\(i=1\), although the chosen basis begins with \(e_0\) and
the ensuing identity sums from \(0\). The lower bound is
therefore \(i=0\) here.}
\endgroup
\phantomsection\label{form:sga3-VIB-11.16-rho-matrix}%
\[
  \rho(e_j)=\sum_{i=0}^{n}e_i\otimes b_{ij}.
\]
As in the proof of
\hyperref[lem:sga3-VIB-11.12]{11.12}, one has
\phantomsection\label{form:sga3-VIB-11.16-star}%
\[
  \tag{\(*\)}
  \Delta(b_{ij})
    =\sum_{\ell=0}^{n}b_{i\ell}\otimes b_{\ell j}.
\]

Put \(a_i=b_{i0}\). Equivalently, if
\((e_0^*,\ldots,e_n^*)\) is the dual basis, then the \(a_i\)
are the matrix coefficients
\[
  g\longmapsto e_i^*(ge_0).
\]
The action of \(H\) on \(E\) corresponds to
\phantomsection\label{form:sga3-VIB-11.16-H-coaction-E}%
\[
  \overline{\rho}:\quad
  E\xrightarrow{\ \rho\ }E\otimes A
    \xrightarrow{\ \mathrm{id}_E\otimes\pi\ }E\otimes B.
\]
Since \(W\) is stable under \(H\), the vector \(e_0\) is
semi-invariant under \(H\). Hence
\[
  \overline{\rho}(e_0)=e_0\otimes\pi(a_0),
  \qquad
  a_i=b_{i0}\in I\quad(1\leq i\leq n).
\]
Substitution in
\hyperref[form:sga3-VIB-11.16-star]{\textup{(\(*\))}} gives
\[
  \overline{\Delta}(a_i)
    =a_i\otimes\pi(a_0)
  \qquad(0\leq i\leq n).
\]
Thus the \(a_i\) are semi-invariant under \(H\), all of
weight \(\pi(a_0)\). Moreover, after replacing \(E\), if
necessary, by the \(G\)-submodule generated by \(e_0\)
\textup{(cf.\ \hyperref[lem:sga3-VIB-11.8]{11.8})}, we may
suppose that the \(a_i\) are linearly independent.

Conversely, let
\[
  H'=\operatorname{Spec}B',
  \qquad
  B'=A/I',
\]
be a closed group subscheme of \(G\) under which every
\(a_i\) is semi-invariant, of weight
\(\lambda_i\in B'\). This includes, in particular, the
case in which \(e_0\) is semi-invariant under \(H'\) of
weight \(\lambda'\). We shall prove that \(H'\subset H\).%
\begingroup
\renewcommand{\thefootnote}{\fnsymbol{footnote}}%
\footnote[3]{\label{note:sga3-VIB-11.16-maximality-source-reading}
Translator's note: the French re-edition says at this
point ``let us show that \(H'=H\),'' as does Reinhold's English translation. The maximality assertion requires
\(H'\subset H\), and that is exactly the conclusion proved
below.}
\endgroup\ 
Let
\(\pi':A\to B'\) be the projection. The hypothesis gives
\phantomsection\label{form:sga3-VIB-11.16-converse-coefficients}%
\[
  a_i\otimes\lambda_i
  =
  (\mathrm{id}_A\otimes\pi')\Delta(b_{i0})
  =
  \sum_{\ell=0}^{n}b_{i\ell}\otimes\pi'(a_\ell).
\]
It follows that
\[
  \lambda_i=\pi'(a_0),
  \qquad
  a_\ell\in I'\quad(1\leq\ell\leq n),
\]
and hence \(e_0\) is semi-invariant under \(H'\). By
\hyperref[lem:sga3-VIB-11.16.0]{Lemma 11.16.0}, \(W\) is
stable under \(H'\). Since \(W\) generates \(I\), the ideal
\(I\) is also stable under \(H'\), and therefore
\[
  \Delta(I)\subset I\otimes A+A\otimes I'.
\]
Finally,
\[
  I\subset\operatorname{Ker}\varepsilon,
  \qquad
  (\varepsilon\otimes\mathrm{id}_A)\circ\Delta
    =\mathrm{id}_A,
\]
so \(I\subset I'\), whence \(H'\subset H\). This completes
the proof.
